\exercisesection{Exercises for § 2}

\begin{enumerate}
\def\labelenumi{\arabic{enumi})}
\item
  Let G be a group, B and N two subgroups of G and S a subset of \(W = N/(B \cap N)\) . For all \(w \in W\) , put \(C(w) = BwB\) . Assume that conditions (T1) and (T2) of Def. 1 of no. 1 are satisfied, that whenever \(s \in S\) and \(w \in W\) at least one of the two relations \(C(s).C(w) = C(sw)\) and \(C(s).C(sw) = C(w)\) is satisfied, and that \(B \cup C(s)\) is a subgroup of G for all \(s \in S\) . Show that condition (T3) is satisfied. Moreover, if the index of B in \(B \cup C(s)\) is \(\geqslant 3\) , then (G, B, N, S) is a Tits system.
\item
  Let G be a group, B and N two subgroups of G and S a subset of \(W = N/(B \cap N)\) . Let Z be a normal subgroup of G contained in B. Let \(B'\) and \(N'\) be the canonical images of B and N in \(G' = G/Z\) . Show that the canonical map from N to \(N'\) defines an isomorphism from W to \(W' = N'/(B' \cap N')\) . Let \(S'\) be the image of S under this isomorphism. Show that \((G', B', N', S')\) is a Tits system if and only if \((G, B, N, S)\) is one.
\item
  Let G be a group, B a subgroup of G and \((\mathrm{C}(w))_{w\in\mathbb{W}}\) the family of double cosets of G with respect to B. Then B is called a Tits subgroup of G if there exists a subset S of W such that the following conditions are satisfied:
\end{enumerate}

\begin{enumerate}
\def\labelenumi{(\arabic{enumi})}
\item
  the union of the \(C(s)\) for \(s \in S\) generates \(G\) ;
\item
  for all \(s \in S\) , the set \(B \cup C(s)\) is a subgroup of \(G\) and the index of \(B\) in \(B \cup C(s)\) is \(\geqslant 3\) ;
\item
  for all \(s \in S\) and all \(w \in W\) , there exists an element \(w' \in W\) such that \(C(s).C(w) \subset C(w) \cup C(w')\) .
\end{enumerate}

From now on we assume that B is a Tits subgroup of G and that we are given a subset S of W satisfying conditions (1), (2) and (3).

\begin{enumerate}
\def\labelenumi{\alph{enumi})}
\item
  Show that \(C(s)^{-1} = C(s)\) and that \(C(s).C(s) = B \cup C(s)\) for all \(s \in S\) . Show that, for all \(s \in S\) and all \(w \in W\) , there exists an element \(w'' \in W\) such that \(C(w).C(s) \subset C(w) \cup C(w'')\) .
\item
  If \(w \in \mathbf{W}\) , the length of \(w\) , denoted by \(l(w)\) , is the smallest integer \(n \geqslant 0\) for which there exist \(s_1, \ldots, s_n \in \mathbf{S}\) with \(\mathrm{C}(w) \subset \mathrm{C}(s_1) \ldots \mathrm{C}(s_n)\) . Show that \(l(w)\) is finite for all \(w \in \mathbf{W}\) .
\end{enumerate}

Let \(u, v \in W\) with \(l(u) < l(v)\) and let \(s \in S\) . Show that, if \(C(v) \subset C(u).C(s)\) (resp. \(C(v) \subset C(s).C(u)\) ), then \(C(v) = C(u).C(s)\) (resp., \(C(v) = C(s).C(u)\) ) (argue by induction on the length of \(u\) . If \(C(v) \neq C(u).C(s)\) , then \(C(u).C(s) = C(u) \cup C(v)\) . By using the induction hypothesis, show that there exist \(t \in S\) and \(w \in W\) such that

\[
\mathrm{C} (u) = \mathrm{C} (t). \mathrm{C} (w) \quad \text { with } \quad l (w) = l (u) - 1.
\]

From the relations \(\mathrm{C}(v) \subset \mathrm{C}(u).\mathrm{C}(s) = \mathrm{C}(t).\mathrm{C}(w).\mathrm{C}(s)\) and \(\mathrm{C}(t).\mathrm{C}(w) = \mathrm{C}(u) \neq \mathrm{C}(v)\) , deduce the existence of an element \(w' \neq w\) such that

\(C(w') \subset C(w).C(s)\) and \(C(v) \subset C(t).C(w')\) , so \(l(w') \geqslant l(v) - 1 > l(u) - 1 = l(w)\) . The induction hypothesis now implies that

\[
\mathrm{C} (w ^ {\prime}) = \mathrm{C} (w). \mathrm{C} (s).
\]

Moreover,

\[
\begin{array}{c} \mathrm{C} (t). \mathrm{C} (u). \mathrm{C} (s) = \mathrm{C} (t). \mathrm{C} (t). \mathrm{C} (w). \mathrm{C} (s) = \mathrm{C} (t). \mathrm{C} (w). \mathrm{C} (s) \cup \mathrm{C} (w). \mathrm{C} (s) \\ = \mathrm{C} (u). \mathrm{C} (s) \cup \mathrm{C} (w ^ {\prime}) = \mathrm{C} (u) \cup \mathrm{C} (v) \cup \mathrm{C} (w ^ {\prime}) \end{array}
\]

and also

\[
\mathrm{C} (t). \mathrm{C} (u). \mathrm{C} (s) = \mathrm{C} (t). (\mathrm{C} (u) \cup \mathrm{C} (v)) \supset \mathrm{C} (t). \mathrm{C} (t). \mathrm{C} (w) \supset \mathrm{C} (w),
\]

which is a contradiction since \(w \neq u, v, w'\) .

\begin{enumerate}
\def\labelenumi{\alph{enumi})}
\setcounter{enumi}{2}
\tightlist
\item
  Show that, for all \(w \in \mathbf{W}\) and all \(s \in S\) , there exists a unique element, denoted by \(s.w\) (resp. \(w * s\) ), distinct from \(w\) and such that
\end{enumerate}

\[
\begin{array}{c} \mathrm{C} (s. w) \subset \mathrm{C} (s). \mathrm{C} (w) \subset \mathrm{C} (w) \cup \mathrm{C} (s. w) \\ (\text {resp.} \mathrm{C} (w * s) \subset \mathrm{C} (w). \mathrm{C} (s) \subset \mathrm{C} (w) \cup \mathrm{C} (w * s)). \end{array}
\]

(Show by induction on \(l(w)\) that \(C(s).C(w) \neq C(w)\) . For this, write

\[
\mathrm{C} (w) = \mathrm{C} (u). \mathrm{C} (t)
\]

with \(t \in S\) , \(u \in W\) and \(l(w) = l(u) + 1\) . If \(C(s).C(w) = C(w)\) , then

\[
\mathrm{C} (s). \mathrm{C} (u). \mathrm{C} (t) = \mathrm{C} (u). \mathrm{C} (t)
\]

and multiplying on the right by \(\mathrm{C}(t)\) , we obtain that \(\mathrm{C}(u) \cup \mathrm{C}(w) = \mathrm{C}(s).\mathrm{C}(u) \cup \mathrm{C}(w)\) . Since \(\mathrm{C}(u) \neq \mathrm{C}(s).\mathrm{C}(u)\) by the induction hypothesis, we have, by b),

\[
\mathrm{C} (s). \mathrm{C} (u) = \mathrm{C} (w),
\]

and hence

\[
\mathrm{C} (u) \subset \mathrm{C} (s). \mathrm{C} (s). \mathrm{C} (u) = \mathrm{C} (s). \mathrm{C} (w) = \mathrm{C} (w),
\]

which is absurd).

\begin{enumerate}
\def\labelenumi{\alph{enumi})}
\setcounter{enumi}{3}
\item
  Let \(s \in S\) . Show that the map \(p_s: w \mapsto s.w\) (resp. \(q_s: w \mapsto w * s\) ) is a permutation of \(W\) and that \(p_s^2 = \operatorname{Id}(\operatorname{resp.} q_s^2 = \operatorname{Id})\) . Show that \(p_s \circ q_t = q_t \circ p_s\) for all \(s, t \in S\) (study the product \(C(s).C(w).C(t)\) for \(w \in W\) and show that \((s.w) * t \in \{w, s.w, w * t, s.(w * t)\}\) ; show that \((s.w) * t \notin \{s.w, w * t\}\) and that if \((s.w) * t = w\) then \(s.w = w * t\) and \(w = s.(w * t)\) ).
\item
  Show that the group of permutations P (resp. Q) generated by the \(p_{s}\) (resp. \(q_{s}\) ) for \(s \in S\) acts simply-transitively on W (to show that P is transitive, use (2) and argue as in Lemma 1 of no. 1; to show that P is simply transitive, use d)). Deduce that W has a unique group structure such that the map \(p \mapsto p(e)\)
\end{enumerate}

(where \(e\) denotes the element of \(W\) such that \(B = C(e)\) ) is an isomorphism from \(P\) to \(W\) . The map \(q \mapsto q(e)\) is then also an isomorphism; moreover, \(s.w = w * s = sw\) for all \(s \in S\) and all \(w \in W\) and \(C(w)^{-1} = C(w^{-1})\) .

\begin{enumerate}
\def\labelenumi{\alph{enumi})}
\setcounter{enumi}{5}
\item
  Show that the pair \((W,S)\) is a Coxeter system and generalize the results of no. 4.
\item
  Let X be a subset of S and \(W_{X}\) the subgroup of W generated by X. Show that the union \(G_{X}\) of the \(C(w)\) for \(w \in W_{X}\) is a subgroup of G and that Th. 3 of no. 5 is still true. Show that B is a Tits subgroup of \(G_{X}\) . Generalize Props. 2 and 3 of no. 5, and Def. 2, Prop. 4 and Th. 4 of no. 6. Show that S is the set of \(w \in W\) such that \(B \cup C(w)\) is a subgroup of G distinct from B. The Coxeter system (W, S) and the group W (called the Weyl group of (G, B)) thus depend only on (G, B).
\item
  Let N be a subgroup of G such that \(B \cap N\) is normal in N and such that the intersection with N of every double coset \(C(w)\) with respect to B is a double coset with respect to \(B \cap N\) . Show that the group \(N/(B \cap N)\) can be identified with W and that \((G, B, N, S)\) is a Tits system.
\end{enumerate}

\begin{enumerate}
\def\labelenumi{\arabic{enumi})}
\setcounter{enumi}{3}
\tightlist
\item
  Let \((G, B, N, S)\) and \((G', B', N', S')\) be two Tits systems with \(G = G'\) and \(B = B'\) , and with Weyl groups W and \(W'\) . Let j be the bijection from W to \(W'\) defined by the relation
\end{enumerate}

\[
\mathrm{B} w \mathrm{B} = \mathrm{B} ^ {\prime} j (w) \mathrm{B} ^ {\prime}.
\]

Show that \(j\) is a group isomorphism between \(W\) and \(W'\) and that \(j(S) = S'\) .

\begin{enumerate}
\def\labelenumi{\arabic{enumi})}
\setcounter{enumi}{4}
\tightlist
\item
  Let \(\Sigma = (G, B, N, S)\) be a Tits system. Put \(T = B \cap N\) and let \(\tilde{N}\) be the normaliser of \(N\) .
\end{enumerate}

\begin{enumerate}
\def\labelenumi{\alph{enumi})}
\tightlist
\item
  Let \(b \in B \cap \hat{N}\) . Show that \(bnb^{-1}n^{-1} \in B \cap N\) for all \(n \in N\) (put \(bn = n'b\) and use Th. 1) and that b belongs to the intersection \(\tilde{T}\) of the conjugates \(nBn^{-1}\) for \(n \in N\) . Show that \(\tilde{T} \cap N = T\) .
\end{enumerate}

If \(\tilde{\mathrm{T}} = \mathrm{T}\) , \(\Sigma\) is said to be saturated.

\begin{enumerate}
\def\labelenumi{\alph{enumi})}
\setcounter{enumi}{1}
\item
  Put \(\tilde{N} = N.\tilde{T}\) . Show that \(\tilde{N}\) is a subgroup of G containing \(\tilde{T}\) as a normal subgroup and that \(\tilde{N} \cap B = \tilde{T}\) . Show that the injection of N into \(\tilde{N}\) defines an isomorphism from the Weyl group W of \(\Sigma\) to \(\tilde{N}/\tilde{T}\) .
\item
  Show that \((\mathbf{G},\mathbf{B},\tilde{\mathbf{N}},j(\mathbf{S}))\) is a saturated Tits system, which is said to be associated to \(\Sigma\) .
\end{enumerate}

\begin{enumerate}
\def\labelenumi{\arabic{enumi})}
\setcounter{enumi}{5}
\item
  We use the notation of no. 2 and let \(\mathbf{N}_0\) be the subgroup of \(\mathbf{N}\) consisting of the matrices all of whose entries are equal to 0 or 1. Show that \(\mathbf{B} \cap \mathbf{N}_0 = \mathbf{T} \cap \mathbf{N}_0 = \{1\}\) and that the canonical map \(j\) from \(\mathbf{N}_0\) to \(\mathbf{W} = \mathbf{N} / \mathbf{T}\) is an isomorphism. Put \(S_0 = j^{-1}(S)\) . Show that \((\mathbf{G}, \mathbf{B}, \mathbf{N}_0, S_0)\) is a Tits system and that \((\mathbf{G}, \mathbf{B}, \mathbf{N}, S)\) is the associated saturated Tits system.
\item
  Let G be a group acting on a set E. Then G is said to act doubly-transitively on E if, whenever \(x, y, x', y' \in E\) are such that \(x \neq y\) and \(x' \neq y'\) , there exists an element \(g \in G\) such that \(g.x = x'\) and \(g.y = y'\) .
\end{enumerate}

\begin{enumerate}
\def\labelenumi{\alph{enumi})}
\item
  Let \((G, B, N, S)\) be a Tits system whose Weyl group W is of order 2. Show that G acts doubly-transitively on G/B.
\item
  Let G be a group acting doubly transitively on a set E. Assume that \(\operatorname{Card}(E) \geqslant 3\) . Let \(e \in E\) and let B be the stabilizer of e. Let \(x \in E\) , with \(x \neq e\) , and let \(n \in G\) be such that \(n(e) = x\) and \(n(x) = e\) . Let N be the subgroup of G generated by n.~Let s be the canonical image of n in N/T. Show that \((G, B, N, \{s\})\) is a Tits system whose Weyl group is of order 2.
\end{enumerate}

\begin{enumerate}
\def\labelenumi{\arabic{enumi})}
\setcounter{enumi}{7}
\tightlist
\item
  Let \((\mathrm{G},\mathrm{B},\mathrm{N},\mathrm{S})\) be a Tits system; put \(T=B\cap N\) and W=N/T. Let \(\tilde{G}\) be a group containing G as a normal subgroup. Assume that for all \(h\in\tilde{G}\) , there exists \(g\in G\) such that \(hBh^{-1}=gBg^{-1}\) and \(hNh^{-1}=gNg^{-1}\) . Let \(\hat{B}\) (resp. \(\hat{N}\) ) be the normaliser of B (resp. N) in \(\tilde{G}\) ; put \(\Gamma=\hat{B}\cap\hat{N}\) , \(\hat{N}=\Gamma.N\) and \(\tilde{T}=\tilde{N}\cap B\) .
\end{enumerate}

\begin{enumerate}
\def\labelenumi{\alph{enumi})}
\item
  Show that \(\tilde{\mathbf{G}} = \Gamma .\mathbf{G}\) , \(\hat{\mathbf{B}} = \Gamma .\mathbf{B}\) , \(\Gamma \cap \mathrm{B} = \Gamma \cap \mathrm{G}\) and \(\tilde{\mathbf{T}} = (\Gamma \cap \mathbf{B}).\mathbf{T}\) . The groups \(\Omega = \Gamma /(\Gamma \cap \mathrm{B})\) , \(\tilde{\mathbf{G}} /\mathbf{G}\) and \(\hat{\mathbf{B}} /\mathbf{B}\) are thus canonically isomorphic. If \(\varPhi \subset \varOmega\) and if \(\mathbf{H}\) is a subgroup of \(\mathbf{G}\) containing \(\Gamma \cap \mathrm{B}\) , we denote by \(\varPhi \mathrm{H}\) the union of the subsets \(\varphi \mathrm{H}\) for \(\varphi \in \varPhi\) .
\item
  Show that \(\tilde{T}\) is normal in \(\tilde{N}\) (to show that \(n\gamma n^{-1} \in \tilde{N}\) for \(n \in \tilde{N}\) and \(\gamma \in \Gamma \cap B\) , use Exerc. 5 a)), that \(N \cap \tilde{T} = T\) and that \(\Gamma \cap \tilde{T} = \Gamma \cap B\) . The injection of N (resp. \(\Gamma\) ) into \(\tilde{N}\) thus allows W (resp. \(\Omega\) ) to be identified with a subgroup of \(\tilde{W} = \tilde{N}/\tilde{T}\) . Show that \(\Omega\) normalises S and that \(\tilde{W}\) is the semi-direct product of \(\Omega\) and W.
\item
  Show that, for all \(s \in S\) and all \(u \in \tilde{W}\) ,
\end{enumerate}

\[
\mathrm{BsBuB} \subset (\mathrm{BuB}) \cup (\mathrm{BsuB}).
\]

\begin{enumerate}
\def\labelenumi{\alph{enumi})}
\setcounter{enumi}{3}
\item
  Show that the map \(u \mapsto BuB\) is a bijection from \(\tilde{W}\) to \(B\backslash\tilde{G}/B\) (use Th. 1 and the fact that \(\Gamma\) normalises B).
\item
  Let \(\mathfrak{B}\) be the set of pairs \((\varPhi, X)\) , where \(\varPhi\) is a subgroup of \(\Omega\) and \(X\) is a subset of \(X\) normalised by \(\varPhi\) . Put \(G_{(\varPhi, X)} = \varPhi G_X = B \varPhi W_X B\) (with the notation of no. 5). Show that the map \((\varPhi, X) \mapsto G_{(\varPhi, X)}\) is a bijection from \(\mathfrak{B}\) to the set of subgroups of \(\tilde{G}\) containing \(B\) . Generalize assertions \(b)\) and \(c)\) of Th. 3 and Prop. 2 of no. 5.
\end{enumerate}

Show that the normaliser of \(\mathrm{G}_{(\varPhi,\mathrm{X})}\) in \(\tilde{G}\) is the subgroup \(\mathrm{G}_{(\varPhi',\mathrm{X})}\) , where \(\varPhi'\) is the set of elements of \(\Omega\) normalizing both \(\varPhi\) and X.

\begin{enumerate}
\def\labelenumi{\alph{enumi})}
\setcounter{enumi}{5}
\tightlist
\item
  Show that \(G_{(\Phi, \mathsf{X})}\) is a maximal subgroup of \(\tilde{\mathbf{G}}\) if and only if one of the following two conditions is satisfied:
\end{enumerate}

\begin{enumerate}
\def\labelenumi{(\roman{enumi})}
\item
  \(X = S\) and \(\varPhi\) is a maximal subgroup of \(\varOmega\) ;
\item
  \(\Phi = \Omega\) and \(\Phi\) acts transitively on \(S - X\) (which is non-empty).
\end{enumerate}

Show that \(G_{(\Phi,X)}\) is a maximal subgroup in the set of subgroups of \(\tilde{G}\) not containing G if and only if

\begin{enumerate}
\def\labelenumi{(\roman{enumi})}
\setcounter{enumi}{2}
\tightlist
\item
  \(X \neq S\) , \(\Phi\) is the normaliser of \(X\) in \(\Omega\) and acts transitively on \(S - X\) .
\end{enumerate}

\(g)\) Let \(\Phi\) be a normal subgroup of \(\Omega\) and put \(G' = \Phi G\) , \(B' = \Phi B\) , \(N' = \Phi N\) and \(T' = B' \cap N'\) . Show that \(T' = \Phi \tilde{T}\) and that \(T'\) is normal in \(N'\) if and only if every element of \(\Phi\) commutes with every element of \(W\) . Show that \(G'\) is then normal in \(\tilde{G}\) , that the injection of \(N\) into \(N'\) defines an isomorphism \(j\) from \(W\) to \(W' = N'/T'\) and that \((G', B', N', S')\) (with \(S' = j(S)\) ) is a Tits system.

\begin{enumerate}
\def\labelenumi{\arabic{enumi})}
\setcounter{enumi}{8}
\tightlist
\item
  Let \((G, B, N, S)\) be a Tits system and X, Y, Z three subsets of S. Show that
\end{enumerate}

\[
\mathrm{G} _ {\mathrm{X}} \cap (\mathrm{G} _ {\mathrm{Y}}. \mathrm{G} _ {\mathrm{Z}}) = (\mathrm{G} _ {\mathrm{X}} \cap \mathrm{G} _ {\mathrm{Y}}). (\mathrm{G} _ {\mathrm{X}} \cap \mathrm{G} _ {\mathrm{Z}})
\]

(use Exerc. 1 of §1 and Prop. 2 of no. 3).

\begin{enumerate}
\def\labelenumi{\arabic{enumi})}
\setcounter{enumi}{9}
\tightlist
\item
  Let G be a group and B a Tits subgroup of G. We use the notation of Exerc. 3. For \(s \in S\) , denote by \(G^{(s)}\) the subgroup \(G_{S - \{s\}}\) (Exerc. 3 g)). Let I be the set of subsets of G of the form \(gG^{(s)}\) (for \(g \in G\) and \(s \in S\) ), and C the set of subsets of I of the form \(C_g = \{gG^{(s)} \mid s \in S\}\) for \(g \in G\) . The \(C_g\) are called the chambers of I (§ 1, Exerc. 15). Let G act on I by left translations.
\end{enumerate}

\begin{enumerate}
\def\labelenumi{\alph{enumi})}
\item
  Let \(\mathfrak{F}\) be the set of facets of I (that is, the subsets of the chambers, cf.~§1, Exerc. 15). Let \(F \in \mathfrak{F}\) . Show that there exist a unique subset X of S and an element \(g \in G\) such that \(gG_{X} = \bigcap_{a \in F} a\) ; F is said to be of type X. Show that F is then the set of \(gG^{(s)}\) for \(s \in S - X\) and is of codimension \(\operatorname{Card}(X)\) in any chamber containing it. Show that the map \(j: F \mapsto \bigcap_{a \in F} a\) is a strictly decreasing bijection, compatible with left translations, from \(\mathfrak{F}\) to the set of subsets of G of the form \(gG_{X}\) for \(g \in G\) and \(X \subset S\) . Show that, if \(X \subset Y \subset S\) , every facet of type X contains a unique facet of type Y.
\item
  Show that G acts transitively on the set C of chambers and that the stabilizer of the chamber \(C_{g}\) ( \(g \in G\) ) is equal to \(gBg^{-1}\) ; the map \(g \mapsto C_{g}\) thus defines a bijection from G/B to C.
\item
  Show that two chambers \(C_{g}\) and \(C_{g'}\) ( \(g, g' \in G\) ) are neighbouring (§1, Exerc. 15) if and only if there exists \(s \in S\) such that \(g' \in g(B \cup BsB)\) .
\end{enumerate}

\(d)\) Let \(C_1, \ldots, C_n\) be chambers of I and put \(C_0 = C_e\) . Show that the following conditions are equivalent:

\begin{enumerate}
\def\labelenumi{(\roman{enumi})}
\item
  the sequence \(\Gamma = (\mathrm{C}_0, \mathrm{C}_1, \ldots, \mathrm{C}_n)\) is an injective gallery;
\item
  there exist a sequence \(\mathbf{s} = (s_1, \ldots, s_n)\) of elements of S and a sequence \((b_1, \ldots, b_n)\) of elements of B such that \(C_j = b_1 \bar{s}_1 b_2 \bar{s}_2 \ldots b_j \bar{s}_j (C_0)\) (where \(\bar{s}_j\) denotes a given element of the double coset \(Bs_jB\) ) for \(1 \leqslant j \leqslant n\) .
\end{enumerate}

Show that, if these conditions are satisfied, the sequence s is unique: it is then called the type of \(\Gamma\) and denoted by \(\mathbf{s}(\Gamma)\) . Show that an injective gallery is minimal if and only if its type is a reduced decomposition. Show that the following conditions are satisfied:

(WI 1) For any two chambers C and \(\mathbf{C}'\) of I, there exists a unique element \(t(\mathbf{C},\mathbf{C}')\) of W such that the set of types of minimal galleries with extremities C and \(\mathbf{C}'\) is the set of reduced decompositions of \(t(\mathbf{C},\mathbf{C}')\) .

(WI 2) For any chamber C, the map \(\mathrm{C}' \mapsto t(\mathrm{C}, \mathrm{C}')\) from the set of chambers of I to W is surjective.

\begin{enumerate}
\def\labelenumi{\alph{enumi})}
\setcounter{enumi}{4}
\item
  Show that two minimal galleries of the same type and with the same extremities are identical. (Reduce to proving that, if \((s_{1},\ldots,s_{n})\) is a reduced decomposition and if \(b_{1},\ldots,b_{n},b_{1}^{\prime},\ldots,b_{n}^{\prime}\) are elements of B such that \(b_{1}\bar{s}_{1}\ldots b_{n}\bar{s}_{n}\in b_{1}^{\prime}\bar{s}_{1}\ldots b_{n}^{\prime}\bar{s}_{n}B\) , then \(b_{1}\bar{s}_{1}\in b_{1}^{\prime}\bar{s}_{1}B\) . For this, remark that if \(\bar{s}_{1}^{-1}b_{1}^{-1}b_{1}^{\prime}\bar{s}_{1}\notin B\) , then this element would belong to \(Bs_{1}B\) and we would have \(b_{2}\bar{s}_{2}\ldots b_{n}\bar{s}_{n}\in b\bar{s}_{1}b^{\prime}b_{2}^{\prime}\bar{s}_{2}\ldots b_{n}^{\prime}\bar{s}_{n}B\) with \(b,b^{\prime}\in B\) , contradicting Cor. 1 of Th. 2 of no. 4.)
\item
  Show that I equipped with C is a building, said to be associated to the pair (G, B). Show that there exists a unique numbering (§ 1, Exerc. 20) of I for which the type of a facet is that defined in a).
\item
  Show that I is spacious, namely that every panel of I is contained in at least three chambers (cf.~§ 1, Exerc. 24). Show that the following condition is satisfied:
\end{enumerate}

\begin{enumerate}
\def\labelenumi{(\Alph{enumi})}
\setcounter{enumi}{6}
\tightlist
\item
  Given a panel F, a chamber C, a gallery \(\Gamma = (\mathrm{C}_0, \ldots, \mathrm{C}_n)\) such that \(\mathrm{F} \subset \mathrm{C}_n\) and of smallest possible length, and chambers \(\mathrm{C}'\) and \(\mathrm{C}''\) containing F and distinct from \(\mathrm{C}_n\) , there exists an element \(g \in \mathrm{G}\) such that \(g(\mathrm{C}_i) = \mathrm{C}_i\) for \(0 \leqslant i \leqslant n\) and \(g(\mathrm{C}') = \mathrm{C}''\) .
\end{enumerate}

(Reduce to the case \(C_{0} = C\) ; let \(u \in S\) be such that F is of type \(S - \{u\}\) and let s be the type of \(\Gamma\) . Show that \((\mathbf{s}, u)\) is a reduced decomposition. Take \(h \in G\) such that \(C_{n} = h(\mathrm{C})\) ; there exist \(b'\) and \(b'' \in B\) such that \(C' = hb'u(\mathrm{C})\) and \(C'' = hb''u(\mathrm{C})\) . Then use Cor. 1 of Th. 2 of no. 4 generalized to the case of a Tits subgroup (cf.~Exerc. 3 f)) to show that there exists \(b \in B\) such that \(bhb'uB = hb''uB\) ; then \(b \in hBuBu^{-1}Bh^{-1} \subset (hBh^{-1}) \cup (hBuBh^{-1})\) . If \(b \in hBuBh^{-1}\) , we would have BhB \(\subset BhBuB = BhuB\) (loc. cit.), which is impossible. Hence \(b(C_{n}) = C_{n}\) and e) implies that \(b(C_{i}) = C_{i}\) for all i.)

\begin{enumerate}
\def\labelenumi{\arabic{enumi})}
\setcounter{enumi}{10}
\tightlist
\item
  Let \((\mathbf{W},\mathbf{S})\) be a Coxeter system and I a building numbered by S (§ 1, Exerc. 20). If \(\Gamma = (\mathbf{C}_0,\dots ,\mathbf{C}_n)\) is an injective gallery, the type of \(\Gamma\) , denoted by \(s(\Gamma)\) , is the sequence \((s_1,\ldots ,s_n)\) of elements of S such that the panel \(\mathbf{C}_{i - 1}\cap \mathbf{C}_i\) is of type \(\mathbf{S} - \{s_i\}\) (for \(1\leqslant i\leqslant n\) ). If conditions (WI 1) and (WI 2) of Exerc. 10 are satisfied, I is called a \((\mathbf{W},\mathbf{S})\) -building.
\end{enumerate}

Let I be a spacious (W, S)-building (cf.~Exerc. 10 g)) and let G be a group of allowed automorphisms of I satisfying the following condition:

(G0) Given three distinct chambers C, C' and C'' containing the same panel, there exists an element \(g \in G\) such that \(g(\mathrm{C}) = \mathrm{C}\) and \(g(\mathrm{C}') = \mathrm{C}'\) .

Choose a chamber C of I and denote by B the stabilizer of C in G.

\begin{enumerate}
\def\labelenumi{\alph{enumi})}
\item
  Show that G acts transitively on the set of chambers of I.
\item
  Let \(C'\) and \(C''\) be two chambers. Show that \(t(\mathrm{C},\mathrm{C}') = t(\mathrm{C},\mathrm{C}'')\) if and only if there exists \(b \in B\) such that \(\mathrm{C}'' = b(\mathrm{C}')\) (if \(t(\mathrm{C},\mathrm{C}') = t(\mathrm{C},\mathrm{C}'') = w\) , consider a reduced decomposition s of w and a minimal gallery \(\Gamma'\) (resp. \(\Gamma''\) ) with extremities C and \(C'\) (resp. \(C''\) ) and of type s. Argue by induction on \(l(w)\) using \((\mathrm{G}_{0})\) and a)). Deduce that there exists a bijection \(w \mapsto \mathrm{B}(w)\) from W to B\textbackslash G/B.
\item
  Let F be the facet of C of type S- \(\{s\}\) . Show that the stabilizer of F in G is equal to B \(\cup\) B(s) (if g(F) = F, then t(C, g(C)) = 1 or s). Show that B \(\subset\) B(s)B(s) (use (G0)).
\item
  Show that B is a Tits subgroup of G and that the Coxeter system of (G,B) is canonically isomorphic to (W,S). (Let \(w \in W\) and \(s \in S\) be such that \(l_{\mathbb{S}}(sw) = l_{\mathbb{S}}(w) + 1\) . Let \(g \in \mathrm{B}(w)\) and \(u \in \mathrm{B}(s)\) ; let \(\mathbf{s} = (s_1, \ldots, s_n)\) be a reduced decomposition of \(w\) , ( \(C = C_0, C_1, \ldots, C_n = g(C)\) ) a minimal gallery of type s and choose an element \(\bar{s}_i \in \mathrm{B}(s_i)\) ; show, by using \((G_0)\) , that there exist elements \(b_i \in B\) such that \(C_i = b_1 \bar{s}_1 \ldots b_i \bar{s}_i(C)\) . Put \(C_i' = u(C_{i-1})\) and show that the gallery \((C, u(C), u(C_1), \ldots, u(C_n))\) is of type \((s, s)\) , and is therefore minimal. Deduce that \(ug \in \mathrm{B}(sw)\) and that \(B(s) B(w) = B(sw)\) . If now \(l(sw) = l(w) - 1\) , put \(w' = sw\) : then \(B(s) B(w') = B(w)\) and hence
\end{enumerate}

\[
\mathrm{B} (s) \mathrm{B} (w) = \mathrm{B} (s) \mathrm{B} (s) \mathrm{B} \left(w ^ {\prime}\right) \subset \mathrm{B} \left(w ^ {\prime}\right) \cup \left(\mathrm{B} (s) \mathrm{B} \left(w ^ {\prime}\right)\right) = \mathrm{B} (s w) \cup \mathrm{B} (w);
\]

finally, since \(\mathrm{B} \subset \mathrm{B}(s)\mathrm{B}(s)\) , we also have \(\mathrm{B}(w') = \mathrm{B}(sw) \subset \mathrm{B}(s)\mathrm{B}(w)\) .

\begin{enumerate}
\def\labelenumi{\alph{enumi})}
\setcounter{enumi}{4}
\tightlist
\item
  Show that there exists a unique isomorphism from the building associated to (G, B) (Exerc. 10) to I, compatible with the action of G and taking the canonical chamber \(C_{e}\) to C.
\end{enumerate}

\begin{enumerate}
\def\labelenumi{\arabic{enumi})}
\setcounter{enumi}{11}
\tightlist
\item
  Let \((\mathbf{G},\mathbf{B},\mathbf{N},\mathbf{S})\) be a Tits system, \(\mathrm{W} = \mathrm{N} / (\mathrm{B}\cap \mathrm{N})\) its Weyl group, I the \((\mathrm{W},\mathrm{S})\) -building associated to \((\mathbf{G},\mathbf{B})\) (Exerc. 10) and \(\mathbf{C} = \mathbf{C}_e\) the canonical chamber of I. Let \(\mathbf{A}_0\) be the apartment associated to the Coxeter system \((\mathrm{W},\mathrm{S})\) (§ 1, Exerc. 16). For all \(g\in \mathbf{G}\) , let \(\varphi_g\) be the map from \(\mathbf{A}_0\) to I that takes a point \(w\mathrm{W}^{(s)}\) of \(\mathbf{A}_0\) ( \(w\in \mathrm{W}, s\in \mathrm{S}\) ) to the point \(gw\mathrm{G}^{(s)}\) of I.
\end{enumerate}

Show that, for all \(g \in G\) , the map \(\varphi_{g}\) is an isomorphism of numbered buildings from \(A_{0}\) to a subset of I that is the union of chambers \(gn(C_{e})\) for \(n \in N\) . Show that I, equipped with the set U of \(\varphi_{g}(A_{0})\) for \(g \in G\) , is a structured building (§1, Exerc. 24) (to prove (SB 2), remark that, if \(g', g'' \in G\) , there exist \(b', b'' \in B\) and \(n \in N\) such that \(g'^{-1}g'' = b'nb''\) ; then put \(g = g'b'n\) and show that \(g'(C)\) and \(g''(C)\) are contained in \(\varphi_{g}(A_{0})\) . To prove (SB 3), reduce to the case of two apartments \(A' = \varphi_{e}(A_{0})\) and

\(A'' = \varphi_b(A_0)\) , with \(b \in B\) , and show that the map \(a \mapsto b(a)\) leaves the points of \(A' \cap A''\) fixed by using Prop. 2 of no. 5). Show that the Coxeter system (W, S) and the numbering of I are adapted to (I, U) (§1, Exerc. 24 e)) and that the set \(\mathfrak{I}\) of allowed isomorphisms from \(A_0\) to the various elements of \(\mathfrak{U}\) is the set of \(\varphi_g\) for \(g \in G\) .

\begin{enumerate}
\def\labelenumi{\arabic{enumi})}
\setcounter{enumi}{12}
\tightlist
\item
  We use the notation of Exerc. 24 of §1: (I, U) is a spacious structured building equipped with a Coxeter system (W, S) and an adapted numbering and I is the set of allowed isomorphisms from the apartment \(A_{0}\) associated to (W, S) to the various apartments of (I, U). Further, let G be a group of allowed automorphisms of I, preserving E: the group G then acts on I and we assume that G acts transitively on I.
\end{enumerate}

Denote by C a chamber of I, A an apartment of U containing C, \(\varphi\) an allowed isomorphism from \(A_{0}\) to A taking the canonical chamber \(C_{e}\) of \(A_{0}\) to C, B the stabilizer of C in G and N the stabilizer of A in G.

\begin{enumerate}
\def\labelenumi{\alph{enumi})}
\item
  Show that, if \(A'\) and \(A''\) are two apartments belonging to U, and containing the same chamber, there exists \(g \in G\) such that \(g(A') = A''\) and that \(g(a) = a\) for all \(a \in A' \cap A''\) . Show that G acts transitively on the set of pairs (A, C) where \(A \in U\) and C is a chamber of A.
\item
  Show that the map \(n \mapsto \varphi^{-1} \circ n \circ \varphi\) is a surjective homomorphism from N to W, with kernel \(B \cap N\) . We can thus identify \(N/(B \cap N)\) with W.
\item
  Show that conditions (WI 1) and (WI 2) of Exerc. 10 d) are satisfied (use the following fact: if an apartment of I contains two chambers C' and C'', it contains every minimal gallery with extremities C' and C'' (§ 1, Exerc. 24 b))).
\item
  Show that condition (G) of Exerc. 10 g) (and a fortiori condition (G₀) of Exerc. 11) is satisfied (with the notation of (G), consider an apartment A' (resp. A'') of I containing C₀ and C' (resp. C''); show by using Exerc. 24 b) of §1 that Γ ⊂ A'' ∩ A'' and use a)).
\item
  Show that \((G, B, N, S)\) is a Tits system and that \((I, \mathfrak{U})\) is canonically isomorphic to the associated numbered structured building (Exerc. 12).
\end{enumerate}

\begin{enumerate}
\def\labelenumi{\arabic{enumi})}
\setcounter{enumi}{13}
\item
  Let \((\mathbf{G},\mathbf{B},\mathbf{N},\mathbf{S})\) be a Tits system and \((\mathbf{I},\mathfrak{U})\) the associated numbered structured building (Exerc. 12). Show that G, considered as a group of automorphisms of I, satisfies the conditions of Exerc. 13. With the notation of Exerc. 12, put \(\mathrm{A}=\varphi_{e}(\mathrm{A}_{0})\) and \(\mathrm{C}=\varphi_{e}(\mathrm{C}_{e})\) ; show that B is the stabilizer of C in G. Let \(\tilde{N}\) be the stabilizer of A in G; show that \(\tilde{\mathrm{N}}/(\mathrm{B}\cap\tilde{\mathrm{N}})\) can be identified with W and that \((\mathrm{G},\mathrm{B},\tilde{\mathrm{N}},\mathrm{S})\) is the saturated Tits system associated to \((\mathrm{G},\mathrm{B},\mathrm{N},\mathrm{S})\) (Exerc. 5).
\item
  We use the hypotheses and notation of Exerc. 10. Assume moreover that the Weyl group \(\mathbf{W}\) of \((\mathbf{G},\mathbf{B})\) is finite and denote by \(w_0\) the longest element of \(\mathbf{W}\) (§1, Exerc. 22). Two chambers \(\mathbf{C}\) and \(\mathbf{C}'\) are said to be opposed if \(t(\mathbf{C},\mathbf{C}') = w_0\) .
\end{enumerate}

\begin{enumerate}
\def\labelenumi{\alph{enumi})}
\item
  Show that if C and C' are opposed, so are C' and C. Show that there exists a chamber opposed to any given chamber C. Show that the stabilizer of a chamber C in G acts transitively on the set of chambers opposed to C.
\item
  Let C and \(C'\) be two opposed chambers. Show that for all \(w \in W\) , there exists a unique chamber \(C_w\) having the following property: if \((s_1, \ldots, s_k)\) (resp. \((s'_1, \ldots, s'_h)\) ) is a reduced decomposition of w (resp. of \(w' = w_0 w^{-1}\) ), there exists a minimal gallery \((C_0 = C, C_1, \ldots, C_n = C')\) of type \((s_1, \ldots, s_k, s'_1, \ldots, s'_h)\) (with \(n = h + k\) ) such that \(C_w = C_k\) (use Exerc. 22 of §1 and Exerc. 10, d) and e)). Show that \(C_w\) and \(C_{ww_0}\) are opposed.
\item
  Let \(\mathfrak{M}\) be the set of pairs of opposed chambers. For \(m = (C, C') \in \mathfrak{M}\) , let \(A_m\) be the union of the chambers \(C_w\) constructed above. Show that I equipped with the set \(\mathfrak{U}\) of \(A_m\) for \(m \in \mathfrak{M}\) is a structured building ( \(\S 1\) , Exerc. 24) and that the Coxeter system (W, S) and the numbering of I are adapted to (I, \(\mathfrak{U}\) ); define a canonical bijection from \(\mathfrak{M}\) to the set denoted by \(\mathfrak{I}\) in Exerc. 24 of \(\S 1\) . We identify these two sets.
\item
  Let \(m = (C, C') \in \mathfrak{M}\) with \(C = C_e\) , and let N be the stabilizer of \(A_m\) in G. Show that \(N/(B \cap N)\) can be identified with W and that \((G, B, N, S)\) is a saturated Tits system (use Exerc. 13).
\end{enumerate}

\begin{enumerate}
\def\labelenumi{\arabic{enumi})}
\setcounter{enumi}{15}
\tightlist
\item
  We retain the hypotheses and notation of Exerc. 15.
\end{enumerate}

\begin{enumerate}
\def\labelenumi{\alph{enumi})}
\item
  Let C and \(C'\) be two chambers. Show that there exists a chamber \(C''\) opposed both to C and to \(C'\) (take \(C''\) opposed to C such that \(t(\mathrm{C},\mathrm{C}')\) is of greatest possible length; if \(t(\mathrm{C}',\mathrm{C}'')\neq w_{0}\) , there exists a neighbouring chamber \(C_{1}\) of \(C''\) such that \(l(t(\mathrm{C}',\mathrm{C}_{1}))>l(t(\mathrm{C}',\mathrm{C}''))\) ; show by using condition (G) of Exerc. 10 that we can assume that \(C_{1}\not\subset A_{(C,C'')}\) and that \(C_{1}\) is then opposed to C).
\item
  Let \(a \in U\) ; show that there exists a unique involutive automorphism \(j_{A}\) (not necessarily allowed) that takes every chamber of A to the opposed chamber (use Exerc. 22 c) of §1). Let F and \(F'\) be two facets of I; show that if \(F' = j_{A}(F)\) for some \(A \in U\) containing F and \(F'\) , then the same is true for all \(A \in U\) containing F and \(F'\) (if \(F, F' \in A \cap A'\) , with \(A, A' \in U\) , consider a chamber C (resp. \(C'\) ) of A (resp. \(A'\) ) containing F (resp. \(F'\) ) and an \(A'' \in U\) containing C and \(C'\) and use (SB 3)). Then F and \(F'\) are said to be opposed. Show that two facets have a common opposed facet if and only if they are of the same type T and that a facet opposed to a facet of type T is of type \(w_{0}Tw_{0}^{-1}\) .
\item
  Let \(A_{0}\) be the apartment associated to a Coxeter system (W, S) (§1, Exerc. 16) and let J be the set of allowed isomorphisms from \(A_{0}\) to the various elements of U. If \(\alpha\) is a half of \(A_{0}\) , with wall L (§1, Exerc. 16) and if \(\varphi \in J\) , we say that \(\varphi(\alpha)\) is a semi-apartment of I, with wall \(\varphi(L)\) . Show that \(\varphi(L)\) then depends only on \(\varphi(\alpha)\) and not on the pair \((\varphi, \alpha)\) . Let \(D_{1}\) and \(D_{2}\) be two distinct semi-apartments, with the same wall L: show that there exist \(\varphi \in J\) and a wall \(L_{0}\) of \(A_{0}\) such that \(L = \varphi(L_{0})\) and that the \(D_{i}\) are the images under \(\varphi\) of the two halves of \(A_{0}\) determined by \(L_{0}\) (choose a panel F contained in L and two chambers \(C_{1}\) and \(C_{2}\) , one contained in \(D_{1}\) and the other in \(D_{2}\) and such that \(C_{1}\) contains F and \(C_{2}\) contains the panel opposed to F in \(D_{1}\) (which is also opposed to F in \(D_{2}\) ) and consider the apartment of I containing \(C_{1}\) and \(C_{2}\) ).
\end{enumerate}

\begin{enumerate}
\def\labelenumi{\arabic{enumi})}
\setcounter{enumi}{16}
\tightlist
\item
  We retain the hypotheses and notation of Exerc. 15 and 16. Choose an element \(\varphi \in \mathfrak{I}\) taking the canonical chamber \(C\) of \(A_0\) ( \(\S 1\) , Exerc. 16) to the canonical chamber \(C_e\) of \(I\) and, as in Exerc. 15 \(d\) ), denote by \(N\) the stabilizer of the apartment \(\varphi(A_0) \in \mathfrak{U}\) in the group \(G\) . For any subset \(D\) of \(A_0\) , denote by \(B_D\) the subgroup of \(G\) leaving fixed all the points of \(\varphi(D)\) : we have \(B = B_C\) and \(B \cap N = B_{A_0}\) .
\end{enumerate}

\begin{enumerate}
\def\labelenumi{\alph{enumi})}
\item
  Let \(\alpha\) be a half of \(A_0\) containing C. Show that \(B_{\alpha}\) acts transitively on the semi-apartments of I distinct from \(\varphi(\alpha)\) whose wall is the wall L of \(\varphi(\alpha)\) (let X be such a semi-apartment, F a panel contained in L and \(C'\) a chamber of \(\varphi(\alpha)\) containing F; show, by using Exerc. 16 c) and Exerc. 13 a), that there exists \(g \in G\) such that \(g(X \cup \varphi(\alpha)) = \varphi(A_0)\) and \(g(C') = C'\) ; show that \(g(F) = F\) and deduce from Exerc. 22 c) of §1 that \(g \in B_{\alpha}\) . Deduce that \(B_L\) acts doubly transitively on the semi-apartments with wall L and that \((B_L, B_{\alpha}, B_L \cap N)\) is a Tits system whose Weyl group is of order 2 (cf.~Exerc. 7).
\item
  Let \(D_{1}\) and \(D_{2}\) be two convex subsets of \(A_{0}\) such that \(C \subset D_{1} \subset D_{2}\) and that there exists a unique half \(\alpha\) of \(A_{0}\) with \(D_{1} \subset \alpha\) and \(D_{2} \not\subset \alpha\) : then \(D_{1} = D_{2} \cap \alpha\) (§1, Exerc. 21 c)). Show that \(B_{D_{1}} = B_{\alpha}B_{D_{2}}\) and that \(B_{\alpha} \cap B_{D_{2}} = B \cap N\) (by considering a gallery of smallest possible length having one extremity equal to C and the other containing a point \(a \in D_{2} - D_{1}\) , show that there exist two neighbouring chambers \(C_{1}'\) and \(C_{2}'\) of \(A_{0}\) such that \(C_{1}' \subset D_{1}\) , \(C_{2}' \subset D_{2}\) , \(C_{2}' \not\subset D_{1}\) and \(C_{1}' \cap C_{2}'\) contained in the wall \(L'\) of \(\alpha\) . Put \(C_{i} = \varphi(C_{i}')\) , \(F = \varphi(C_{1}' \cap C_{2}')\) and \(L = \varphi(L')\) . Show that the convex hull of \(D_{1} \cup C_{2}'\) is equal to \(D_{2}\) . Let \(b \in B_{D_{1}}\) : then \(b(C_{1}) = C_{1}\) , so \(b(F) = F\) and \(b(C_{2}) \neq C_{1}\) . Deduce that the convex hull of \(b(C_{2}) \cup L\) is a semi-apartment X distinct from \(\varphi(\alpha)\) and with wall L, and that there exists \(b' \in B_{\alpha}\) such that
\end{enumerate}

\[
b ^ {\prime} (\varphi (A _ {0})) = X \cup \varphi (\alpha).
\]

Show that \(b'b(a) = a\) for all \(a \in \varphi(\mathrm{D}_1 \cup \mathrm{C}_2')\) and that \(b'b \in \mathrm{B}_{\mathrm{D}_2}\) .

\begin{enumerate}
\def\labelenumi{\alph{enumi})}
\setcounter{enumi}{2}
\tightlist
\item
  Let \((\alpha_{i})_{1\leqslant i\leqslant q}\) be the halves of \(A_0\) containing C, numbered so that the map
\end{enumerate}

\[
j \mapsto \bigcap_ {1 \leqslant i \leqslant j} \alpha_ {i}
\]

is strictly decreasing (cf.~§1, Exerc. 17). Show that \(\mathbf{B} = \mathbf{B}_{\alpha_1} \ldots \mathbf{B}_{\alpha_q}\) and that if \(b_i, b_i' \in \mathbf{B}_{\alpha_i}\) with \(b_1 \ldots b_q = b_1' \ldots b_q'\) , then \(b_i' \in b_i(\mathbf{B} \cap \mathbf{N})\) for \(1 \leqslant i \leqslant q\) .

\begin{enumerate}
\def\labelenumi{(\arabic{enumi})}
\setcounter{enumi}{17}
\tightlist
\item
  We use the hypotheses and notation of no. 2. Let \(V_i\) be the vector subspace of \(k^n\) generated by \(e_1, \ldots, e_i\) (for \(1 \leqslant i \leqslant n - 1\) ).
\end{enumerate}

\begin{enumerate}
\def\labelenumi{\alph{enumi})}
\item
  Show that, for any subset X of S, the subgroup \(G_{X}\) (no. 5) consists of the elements \(g \in G\) such that \(g(V_{i}) = V_{i}\) for all i such that \(s_{i} \notin X\) .
\item
  Let I be the building associated to a Tits system (G, B, N, S) (Exerc. 10 and 12). Show that the map j that associates to the point \(g\mathbf{G}^{(i)}\) of I (where \(\mathbf{G}^{(i)}\) denotes the subgroup \(G_{S-\{s_i\}}\) of G, in other words the stabilizer of \(V_i\) ) the vector subspace \(g(V_i)\) is a bijection from I to the set B of vector subspaces of \(k^n\) that are \(\neq \{0\}\) and \(\neq k^n\) , compatible with the action of G.
\item
  If E is a vector space, a flag of E is a set of vector subspaces of E that is totally ordered by inclusion. Show that elements \(a_{1}, \ldots, a_{k}\) of I belong to the same facet of I if and only if \(\{j(a_{1}), \ldots, j(a_{k})\}\) is a flag of \(k^{n}\) .
\item
  Show that G acts doubly-transitively on the set of 1-dimensional vector subspaces of \(k^{n}\) . Suppose that \(n \geqslant 2\) and let \(N_{1}\) be the subgroup of G generated by the element that interchanges \(e_{1}\) and \(e_{2}\) and leaves fixed the other \(e_{i}\) . Show that \((\mathrm{G}, \mathrm{G}^{(1)}, \mathrm{N}_{1})\) is a Tits system whose Weyl group is of order 2.
\item
  Assume that \(k\) is commutative. Set
\end{enumerate}

\[
\mathrm{G} ^ {\prime} = \mathbf {S L} (n, k), \quad \mathrm{B} ^ {\prime} = \mathrm{G} ^ {\prime} \cap \mathrm{B}, \quad \mathrm{N} ^ {\prime} = \mathrm{G} ^ {\prime} \cap \mathrm{N} \text {and} \mathrm{T} ^ {\prime} = \mathrm{N} ^ {\prime} \cap \mathrm{B} ^ {\prime} = \mathrm{T} \cap \mathrm{G} ^ {\prime}.
\]

Show that \(N^{\prime}/T^{\prime}\) can be identified with \(S_{n}\) and that \((G^{\prime}, B^{\prime}, N^{\prime}, S)\) is a Tits system (argue as in no. 2).

\begin{enumerate}
\def\labelenumi{\arabic{enumi})}
\setcounter{enumi}{18}
\item
  Let \(k\) be a commutative field of characteristic \(\neq 2\) and let \(\mathbf{Q}\) be the quadratic form \(x_{1}x_{3} + x_{2}^{2}\) on \(k^{3}\) . Show that the group \(\mathbf{SO}(\mathbf{Q})\) acts doubly transitively on the set of isotropic lines in \(k^{3}\) . Compare the corresponding Tits system (Exerc. 7) with that obtained in no. 2 for \(n = 2\) (cf.~Algebra Chap. IX, §9, Exerc. 15).
\item
  We use the notation of no. 2, and assume that k is commutative. Consider the following cases:
\end{enumerate}

\((B_{l})\ n=2l+1\ (with\ l\geqslant1),\ k\ is\ of\ characteristic\neq2\ and\ k^{n}\ is\ equipped\) with the quadratic form \(Q=x_{1}x_{n}+\cdots+x_{l}x_{l+2}+x_{l+1}^{2};\)

\((C_{l}) \ n = 2l \ (\text{with } l \geqslant 1) \ \text{and} \ k^{n} \ \text{is equipped with the alternating form } \Phi \ \text{such that } \Phi(e_{i}, e_{j}) = 0 \ \text{for} \ 1 \leqslant i \leqslant n, \ 1 \leqslant j \leqslant n \ \text{and} \ i \leqslant j, \ \text{except when} \ i \leqslant l, \ j = n + 1 - i, \ \text{when} \Phi(e_{i}, e_{j}) = 1;\)

\((D_{l})\ n = 2l \text{ (with } l \geqslant 2), k \text{ is of characteristic } \neq 2 \text{ and } k^{n} \text{ is equipped with the quadratic form } Q = x_{1}x_{n} + \cdots + x_{l}x_{l+1}.\)

Denote by \(G_{1}\) the special orthogonal group \(\mathbf{SO}(Q)\) in cases \((B_{l})\) and \((D_{l})\) , and the symplectic group \(\mathbf{Sp}(\varPhi)\) in case \((C_{l})\) . Put

\[
\mathrm{B} _ {1} = \mathrm{G} _ {1} \cap \mathrm{B}, \quad \mathrm{N} _ {1} = \mathrm{G} _ {1} \cap \mathrm{N} \text {and} \mathrm{T} _ {1} = \mathrm{G} _ {1} \cap \mathrm{T} = \mathrm{B} _ {1} \cap \mathrm{N} _ {1}.
\]

\begin{enumerate}
\def\labelenumi{\alph{enumi})}
\tightlist
\item
  Show that the action of \(N_{1}\) on the set of lines \(ke_{i}\) defines a homomorphism from \(N_{1}\) , with kernel \(T_{1}\) , onto the subgroup \(W_{1}\) of \(S_{n}\) , allowing us to identify \(W_{1}\) with \(N_{1}/T_{1}\) . Show that \(W_{1}\) is the subgroup of \(S_{n}\) generated by
\end{enumerate}

the \(\sigma_{j} = s_{j}s_{n - j}\) for \(1\leqslant j <   l\) and by \(\sigma_l = s_{l - 1}s_ls_{l - 1}\) in case \((\mathrm{B}_l)\) ;

the \(\sigma_{j} = s_{j}s_{n - j}\) for \(1\leqslant j <   l\) and by \(\sigma_l = s_l\) in case \((C_l)\) ;

\[
\text { the } \sigma_ {j} = s _ {j} s _ {n - j} \text { for } 1 \leqslant j <   l \text { and by } \sigma_ {l} = s _ {l - 1} s _ {l} s _ {l - 1} s _ {l} s _ {l + 1} s _ {l} \text { in case } (D _ {l}).
\]

\begin{enumerate}
\def\labelenumi{\alph{enumi})}
\setcounter{enumi}{1}
\tightlist
\item
  Let \(S_{1}\) be the set of \(\sigma_{j}\) for \(1 \leqslant j \leqslant l\) . Show that \((\mathrm{G}_{1}, \mathrm{B}_{1}, \mathrm{N}_{1}, \mathrm{S}_{1})\) is a Tits system (to prove that the subgroup H of \(G_{1}\) generated by \(B_{1}\) and \(N_{1}\) is the whole of \(G_{1}\) , argue as in Algebra, Chap. II, §10, no. 13), remarking that H contains the lower triangular subgroup of \(G_{1}\) and that, for any \(\xi_{i} \in k\) ( \(2 \leqslant i \leqslant n\) ), there exists a matrix \(b = (b_{ij}) \in B_{1}\) such that \(b_{11} = 1\) , \(b_{1i} = \xi_{i}\) for \(2 \leqslant i \leqslant n - 1\) and \(b_{1n} = \xi_{n}\) in case \((\mathrm{C}_{l})\) , and \(b_{1n} = 0\) in cases \((\mathrm{B}_{l})\) and \((\mathrm{D}_{l})\) . Then argue as in no. 2, introducing the subgroups \(G_{1,j} = G_{1} \cap (G_{j}, G_{n-j})\) for \(1 \leqslant j < l\) and the subgroup \(G_{1,l}\) of elements of \(G_{1}\) leaving fixed
\end{enumerate}

the \(e_i\) for \(i \neq l, l + 1, l + 2\) and the subspace generated by \(e_l, e_{l+1}\) and \(e_{l+2}\) in case \((B_l)\) ;

the \(e_{i}\) for \(i \neq l, l + 1\) and the plane generated by \(e_{l}\) and \(e_{l+1}\) in case \((\mathrm{C}_{l})\) ; the \(e_{i}\) for i \textless{} l - 1 or \(i > l + 2\) and the two planes generated by \(e_{l-1}\) and \(e_{l+1}\) and by \(e_{l}\) and \(e_{l+2}\) in case \((\mathrm{D}_{l})\) .

Show that \(G_{1,j}\) can be identified, respectively, with \(\mathbf{GL}(2,k)\) , \(\mathbf{SL}(2,k)\) or the special orthogonal group of Exerc. 19).

\(^{*}c)\) Show that the Coxeter graph of the group \(W_{1}\) is of type \((B_{l})\) , \((C_{l})\) or \((D_{l})\) in the three cases (Chap VI, §4, no. 1).

\(d\) ) Show that, for any subset X of \(\mathbf{S}_1\) , the subgroup \(\mathbf{G}_{1\mathbf{X}}\) consists of the \(g\in \mathbf{G}_1\) such that \(g(\mathrm{V}_i) = \mathrm{V}_i\) for all \(i\) such that \(\sigma_{i}\notin \mathbf{X}\) , except in case \((\mathrm{D}_l)\) where the same assertion is valid provided that \(\mathrm{V}_{r - 1}\) denotes the subspace generated by \(e_1,\dots ,e_{r - 1}\) and \(e_{r + 1}\) . Deduce, as in Exerc. 18 b), that there exists a bijection \(j\) from the building associated to \((\mathrm{G}_1,B_1)\) to the set of totally isotropic subspaces \(\neq 0\) in cases \((\mathrm{B}_l)\) and \((\mathrm{C}_l)\) , and to the set of totally isotropic subspaces of dimension \(\neq 0\) and \(\neq r - 1\) in case \((\mathrm{D}_l)\) . Show that points \(a_1,\ldots ,a_k\) of I belong to the same facet if and only if \(\{j(a_1),\dots ,j(a_k)\}\) is a flag.

\begin{enumerate}
\def\labelenumi{\arabic{enumi})}
\setcounter{enumi}{20}
\tightlist
\item
  Let A be a discrete valuation ring (Commutative Algebra, Chap. VI, §3, no. 6), m its maximal ideal, γ a generator of m and K the field of fractions of A. Let G be the group SL(2,K), B the subgroup of G consisting of the matrices \(\begin{pmatrix} a & b \\ c & d \end{pmatrix}\) such that \(a, b, d \in A\) and \(c \in m\) (with ad - bc = 1) and N the subgroup consisting of the matrices belonging to G and having only one non-zero entry in each row and column.
\end{enumerate}

\begin{enumerate}
\def\labelenumi{\alph{enumi})}
\item
  Show that \(\mathrm{T} = \mathrm{B} \cap \mathrm{N}\) is normal in \(\mathbb{N}\) and that \(\mathrm{W} = \mathrm{N} / \mathrm{T}\) is an infinite dihedral group generated by the classes \(s\) and \(s'\) of the matrices \(\left( \begin{array}{cc}0 & 1\\ -1 & 0 \end{array} \right)\) and \(\left( \begin{array}{cc}0 & \gamma \\ -\gamma^{-1} & 0 \end{array} \right)\) , respectively.
\item
  Show that \((G, B, N, S)\) (with \(S = \{s, s'\}\) ) is a Tits system.
\item
  Let \(\mathbf{H} = \mathbf{SL}(2, \mathbf{A})\) be the subgroup of G consisting of the matrices with coefficients in A. Show that \((\mathbf{H}, \mathbf{B} \cap \mathbf{H}, \mathbf{N}, \{s\})\) is a Tits system. Compare with Exerc. 18 e).
\item
  Let \(\hat{\mathbf{A}}\) be the completion of \(\mathbf{A}\) and let \(\hat{\mathbf{G}},\hat{\mathbf{B}},\hat{\mathbf{N}},\hat{\mathbf{T}}\) be the groups defined as above but replacing \(\mathbf{A}\) by \(\hat{\mathbf{A}}\) . Show that the injection of \(\mathbf{G}\) into \(\hat{\mathbf{G}}\) defines an isomorphism from the building \(\mathbf{I}\) associated to \((\mathbf{G},\mathbf{B})\) (Exerc. 10) to the building \(\hat{\mathbf{I}}\) associated to \((\hat{\mathbf{G}},\hat{\mathbf{B}})\) . Let \((\mathbf{I},\mathfrak{U})\) (resp. \((\hat{\mathbf{I}},\hat{\mathfrak{U}})\) ) be the structured building associated to \((\mathbf{G},\mathbf{B},\mathbf{N})\) (resp. \((\hat{\mathbf{G}},\hat{\mathbf{B}},\hat{\mathbf{N}})\) ) (Exerc. 12): show that \(j(\mathfrak{U})\subset \hat{\mathfrak{U}}\) , but that \(j(\mathfrak{U})\neq \hat{\mathfrak{U}}\) if \(\mathbf{A}\neq \hat{\mathbf{A}}\) (remark that the apartments \(\hat{\mathfrak{U}}\) (resp. \(\mathfrak{U}\) ) correspond bijectively to the conjugates of \(\mathbf{T}\) (resp. \(\hat{\mathbf{T}}\) ) by \(\mathbf{G}\) (resp. \(\hat{\mathbf{G}}\) ).
\end{enumerate}

\begin{enumerate}
\def\labelenumi{\arabic{enumi})}
\setcounter{enumi}{21}
\tightlist
\item
  Let \(G\) be a group and \(B\) a subgroup of \(G\) .
\end{enumerate}

\begin{enumerate}
\def\labelenumi{\alph{enumi})}
\tightlist
\item
  Show that the following conditions are equivalent:
\end{enumerate}

\begin{enumerate}
\def\labelenumi{(\roman{enumi})}
\item
  \(B \cap gBg^{-1}\) is of finite index in \(B\) for all \(g \in G\) ;
\item
  every double coset BgB with respect to B is a finite union of left cosets with respect to B.
\end{enumerate}

More precisely, show that, for all \(g \in G\) , the index \(q_{g}\) of \(B \cap gBg^{-1}\) in B is equal to the number of left cosets with respect to B contained in the double coset BgB. Show that \(q_{gh} \leqslant q_{g}q_{h}\) for all \(g, h \in G\) .

Assume from now on that conditions (i) and (ii) are satisfied and denote by k a commutative ring. For \(t \in G/B\) (resp. \(t \in B \backslash G/B\) ), denote by \(a_{t}\) the map from G to k defined by \(a_{t}(g) = 0\) if \(g \notin t\) and \(a_{t}(g) = 1\) if \(g \in t\) . Let L (resp. H) be the k-module generated by the \(a_{t}\) for \(t \in G/B\) (resp. \(t \in B \backslash G/B\) ).

\begin{enumerate}
\def\labelenumi{\alph{enumi})}
\setcounter{enumi}{1}
\item
  Show that there exists a unique linear form \(\mu\) on \(\mathbf{L}\) such that \(\mu(a_t) = 1\) for all \(t \in \mathbf{G} / \mathbf{B}\) .
\item
  Let \(\varphi \in \mathbf{L}\) and \(\psi \in \mathbf{H}\) . Show that, for all \(x \in \mathbf{G}\) , the map
\end{enumerate}

\[
\theta_ {x}: y \mapsto \varphi (y) \psi (y ^ {- 1} x)
\]

belongs to L and that the map \(\varphi * \psi : x \mapsto \mu(\theta_x)\) belongs to L. Moreover, if \(\varphi \in H\) then \(\varphi * \psi \in H\) . Show that the map \((\varphi, \psi) \mapsto \varphi * \psi\) makes H into an algebra over \(k\) , having \(a_B\) as its unit element, and makes L into a right H-module.

The algebra \(\mathbf{H}\) is called the Hecke algebra of \(\mathbf{G}\) with respect to \(\mathbf{H}\) and is denoted by \(\mathrm{H}_k(\mathbf{G},\mathbf{B})\) .

\begin{enumerate}
\def\labelenumi{\alph{enumi})}
\setcounter{enumi}{3}
\tightlist
\item
  Show that, for \(t, t' \in B \backslash G / B\) ,
\end{enumerate}

\[
a _ {t} * a _ {t ^ {\prime}} = \sum_ {t ^ {\prime \prime}} m (t, t ^ {\prime}; t ^ {\prime \prime}) a _ {t ^ {\prime \prime}},
\]

where \(m(t, t'; t'')\) is the number of cosets with respect to B contained in \(t \cap gt'^{-1}\) for all \(g \in t''\) .

\begin{enumerate}
\def\labelenumi{\alph{enumi})}
\setcounter{enumi}{4}
\tightlist
\item
  Let G act on L by left translations. Show that the action of H on L defines an isomorphism from H to the commutant of the linear representation of G on L thus obtained.
\end{enumerate}

\(^{*}f\) ) Assume that G is finite and that the characteristic of k does not divide the order of G. Show that \(H_{k}(G,B)\) is absolutely semi-simple over k (Algebra, Chap VIII, §7, no. 5) (use Maschke's Theorem (Chap. V, Appendix) and Prop. 3 of Algebra, Chap. VIII, §5, no. 1).

\begin{enumerate}
\def\labelenumi{\alph{enumi})}
\setcounter{enumi}{6}
\tightlist
\item
  Assume that G is a topological group and that B is a compact open subgroup of G. Show that conditions (i) and (ii) are satisfied and that, when k = R or C, the product \(\varphi * \psi\) is simply the convolution product relative to the right Haar measure on G, normalised by the condition \(\mu(B) = 1\) (cf.~Integration, Chap. VIII, §4, no. 5).
\end{enumerate}
% semantic-fragments: legacy-input-seam
\relax{}
\begin{enumerate}
\def\labelenumi{\arabic{enumi})}
\setcounter{enumi}{22}
\tightlist
\item
  Let \((\mathrm{W},\mathrm{S})\) be a Coxeter system and \(k\) a commutative ring. Suppose that we are given, for all \(s\in S\) , two elements \(\lambda_{s}\) and \(\mu_s\) of \(k\) such that \(\lambda_{s} = \lambda_{s'}\) and \(\mu_{s} = \mu_{s'}\) whenever \(s\) and \(s'\) are conjugate in W. Put \(\mathbf{E} = k^{(\mathbf{W})}\) and let \(\{e_w\}\) be the canonical basis of E.
\end{enumerate}

\begin{enumerate}
\def\labelenumi{\alph{enumi})}
\tightlist
\item
  Show that \(E\) has a unique algebra structure such that, for \(s \in S\) and \(w \in W\) ,
\end{enumerate}

\[
e _ {s}. e _ {w} = \left\{ \begin{array}{l l} e _ {s w} & \text {if} l (s w) > l (w) \\ \lambda_ {s} e _ {w} + \mu_ {s} e _ {s w} & \text {if} l (s w) <   l (w) \end{array} \right.
\]

(introduce the endomorphism \(P_{s}\) of E defined by the above formulas, where \(e_{s}.e_{w}\) is replaced by \(\mathrm{P}_{s}(w)\) , and the endomorphism \(Q_{s}=jP_{s}j^{-1}\) , where j denotes the automorphism of E defined by \(j(e_{w})=e_{w^{-1}}\) ; show that \(P_{s}Q_{t}=Q_{t}P_{s}\) for \(s,t\in S\) by remarking that the conditions \(l(swt)=l(w)\) and \(l(sw)=l(wt)\) imply that sw=wt; then argue as in Exerc. 3 c)). The module E, equipped with this algebra structure, will be denoted by \(\mathrm{E}_{k}((\lambda_{s}),(\mu_{s}))\) . Show that \(\mathrm{E}((0),(1))\) is the group algebra k{[}W{]} of W (Algebra, Chap. III, §2, no. 6).

b) Show that the family of generators \((e_s)_{s\in S}\) and the relations

\[
e _ {s} ^ {2} = \lambda_ {s} e _ {s} + \mu_ {s} \mathrm{for} s \in S
\]

\((e_s e_t)^r = (e_t e_s)^r\) for \(s, t \in S\) such that \(st\) is of finite even order \(2r\) \((e_s e_t)^r e_s = (e_t e_s)^r e_t\) for \(s, t \in S\) such that \(st\) is of finite odd order \(2r + 1\)

form a presentation of the algebra \(\mathbf{E}\) (argue as in the proof of Th. 1 of no. 6 of §1).

\begin{enumerate}
\def\labelenumi{\arabic{enumi})}
\setcounter{enumi}{23}
\tightlist
\item
  Let G be a group and B a Tits subgroup of G (Exerc. 3, of which we use the notation). Assume that, for all \(s \in S\) , the double coset \(C(s)\) is the union of a finite number \(q_{s}\) of left cosets with respect to B. We use the notation of Exerc. 22 and put \(a_{w} = a_{\mathrm{C}(w)}\) for all \(w \in W\) .
\end{enumerate}

\begin{enumerate}
\def\labelenumi{\alph{enumi})}
\item
  Show that conditions (i) and (ii) of Exerc. 22 are satisfied. We can therefore speak of the Hecke algebra \(\mathrm{H}_k(\mathbf{G},\mathbf{B})\) ( \(k\) being a commutative ring), of which \((a_w)_{w\in \mathbb{W}}\) is a basis.
\item
  Let \(s \in S\) and \(w \in W\) . Show that \(a_s * a_s = (q_s - 1)a_s + q_s\) ; show that if \(l(sw) > l(w)\) then \(a_s * a_w = a_{sw}\) .
\item
  Show that the linear map from the algebra \(\mathrm{E}_{k}((q_{s}-1),(q_{s}))\) associated to the Coxeter system (W,S) (Exerc. 23) to \(\mathrm{H}_{k}(\mathrm{G},\mathrm{B})\) that sends \(e_{w}\) to \(a_{w}\) for all \(w \in W\) is an isomorphism of algebras.
\end{enumerate}

\begin{enumerate}
\def\labelenumi{\arabic{enumi})}
\setcounter{enumi}{24}
\tightlist
\item
  We use the notation of Exerc. 8. Assume in addition that, for all \(s \in S\) , the index \(q_s\) of \(B \cap gBg^{-1}\) in \(B\) is finite for all \(g \in BsB\) .
\end{enumerate}

\begin{enumerate}
\def\labelenumi{\alph{enumi})}
\item
  Show that the pair \((\tilde{\mathbf{G}},\mathbf{B})\) satisfies conditions (i) and (ii) of Exerc. 22.
\item
  Show that, for all \(\gamma \in \Gamma\) , the map \(x \mapsto \gamma x \gamma^{-1}\) defines an automorphism \(\sigma\) of the Hecke algebra \(\mathrm{H}_k(\mathrm{G}, \mathrm{B})\) and that \(\sigma\) depends only on the class \(\omega\) of \(\gamma\) in \(\Omega = \Gamma / (\Gamma \cap \mathrm{B})\) .
\item
  Let \(k[\Omega]\) be the group algebra of \(\Omega\) and \((e_{\omega})\) its canonical basis. Show that the linear map \(j\) from \(k[\Omega] \otimes_k \mathrm{H}_k(\mathbf{G}, \mathbf{B})\) to \(\mathrm{H}_k(\tilde{\mathrm{G}}, \mathrm{B})\) defined by \(j(e_{\omega} \otimes a_{\mathrm{B}w\mathrm{B}}) = a_{\mathrm{B}\omega w\mathrm{B}}\) (in the notation of Exerc. 22), for \(\omega \in \Omega\) and \(w \in W\) , is bijective and that
\end{enumerate}

\[
j ^ {- 1} (j (e _ {\omega} \otimes x) j (e _ {w ^ {\prime}} \otimes y)) = e _ {w w ^ {\prime}} \otimes \sigma_ {\omega} (x) y
\]

for \(\omega, \omega' \in \Omega\) and \(x, y \in \mathrm{H}_k(\mathrm{G}, \mathrm{B})\) .

\begin{enumerate}
\def\labelenumi{\arabic{enumi})}
\setcounter{enumi}{25}
\tightlist
\item
  If E is an absolutely semi-simple algebra of finite rank over a commutative field k, the numerical invariant of E is the sequence of integers \((n_{1},\ldots,n_{r})\) such that \(n_{1}\geqslant\cdots\geqslant n_{r}>0\) and that \(\bar{k}\otimes_{k}E\) is isomorphic, for any algebraic closure \(\bar{k}\) of k, to \(\prod M_{n_{i}}(\bar{k})\) .
\end{enumerate}

Let V be an integral ring, K its field of fractions, \(\varphi\) a homomorphism from V to a commutative field k, and E a V-algebra. Assume that E is a free V-module of finite rank and put \(E_{0} = E \otimes_{V} k\) and \(E_{1} = E \otimes_{V} K\) .

\begin{enumerate}
\def\labelenumi{\alph{enumi})}
\item
  Assume that the bilinear form \((x,y)\mapsto\mathrm{Tr}_{\mathbb{E}_{0}/k}(xy)\) on \(E_{0}\) is nondegenerate. Show that \(E_{0}\) and \(E_{1}\) are absolutely semi-simple (cf.~Algebra, Chap IX, §2, Exerc. 1).
\item
  Assume that \(E_{0}\) and \(E_{1}\) are absolutely semi-simple over k and K, respectively. Show that \(E_{0}\) and \(E_{1}\) have the same numerical invariant (reduce to the case where k and K are algebraically closed); let \((e_{i})\) be a basis of E over V and let \((\mathbf{X}_i)\) be indeterminates. Show that the characteristic polynomial of the element \(\sum_{i}\mathbf{X}_{i}e_{i}\) of \(\mathrm{E}_1\otimes_{\mathrm{K}}\mathrm{K}[(\mathrm{X}_i)]\) (resp. \(\mathrm{E}_0\otimes_kk[(\mathrm{X}_i)]\) ) is of the form \(\mathrm{P} = \prod_j\mathrm{P}_j^{n_j}\) (resp. \(\mathrm{Q} = \prod_k\mathrm{Q}_k^{m_k}\) ), where \((n_1,\ldots ,n_r)\) (resp. \((m_1,\dots ,m_s))\) is the numerical invariant of \(\mathrm{E}_1\) (resp. \(\mathrm{E}_0\) ), with \(\deg (\mathrm{P}_j) = n_j\) (resp. \(\deg (\mathrm{Q}_k) = m_k\) ). Show that \(\mathrm{P}_j\in \mathrm{V}[(\mathrm{X}_i)]\) and that \(\mathrm{Q} = \varphi (\mathrm{P})\) . Deduce that there exist integers \(c_{jk}\geqslant 0\) such that
\end{enumerate}

\[
m _ {k} = \sum_ {j} c _ {j k} n _ {j}, \text { and } n _ {j} = \sum_ {k} c _ {j k} m _ {k}.
\]

*27) Let \((\mathrm{G},\mathrm{B},\mathrm{N},\mathrm{S})\) be a Tits system and let \(k\) be a commutative field. Assume that G is finite and that the characteristic of \(k\) divides neither the order of G nor the order of the Weyl group \(\mathrm{W} = \mathrm{N} / (\mathrm{B}\cap \mathrm{N})\) . Show that the algebras \(\mathrm{H}_k(\mathrm{G},\mathrm{B})\) (Exerc. 22) and \(k[\mathrm{W}]\) are absolutely semi-simple and have the same numerical invariant, and hence are isomorphic when \(k\) is algebraically closed (let \(q_{s}\) be the index of \(\mathrm{B}\cap g\mathrm{Bg}^{-1}\) in B for \(g\in \mathrm{BsB}\) , \(s\in S\) . Consider the algebra \(\mathrm{E}_{k[\mathrm{X}]}((\mathrm{X}(q_s - 1)),(1 + \mathrm{X}(q_s - 1)))\) constructed as in Exerc. 23 from the Coxeter system (W,S) and the polynomial ring \(k[\mathrm{X}]\) and use Exerc. 26 a) and b), remarking that, by Maschke's theorem (Chap. V, Appendix), the bilinear form \(\mathrm{Tr}_{k[\mathrm{W}] / k}(xy)\) is non-degenerate).*

\begin{enumerate}
\def\labelenumi{\arabic{enumi})}
\setcounter{enumi}{27}
\item
  Let G be a group, M a maximal subgroup of G and U a normal subgroup of M, satisfying condition (R) of no. 7. Assume that G is equal to its commutator subgroup, that it is generated by the union of the conjugates of U and that the intersection of the conjugates of M reduces to the unit element. Show that G is simple. (Consider a normal subgroup N of G distinct from \{1\}; show that G = N.M and then that G = N.U.)
\item
  Let H be a simple non-abelian group. Let \(\theta\) be an automorphism of H of prime order p and let U be the semi-direct product of Z/pZ by H corresponding to \(\theta\) . Show that, if \(\theta\) is not an inner automorphism, the only normal subgroups of U are \(\{1\}\) , H and U. Deduce that U is not soluble, but that it satisfies condition (R) of no. 7. Apply this to show that the symmetric group \(S_{n}\) satisfies condition (R).
\end{enumerate}

